% topic: algorithms
\question กำหนดให้หุ่นยนต์ตัวหนึ่งต้องการเดินทางในเขาวงกตขนาด $6 \times 6$ ช่อง โดยเริ่มเดินจากพิกัดซ้ายบนที่ $(0,0)$ ไปยังพิกัดขวาล่างที่ $(5,5)$ โดยหุ่นยนต์สามารถเดินได้ 4 ทิศทาง (บน, ล่าง, ซ้าย, ขวา) ครั้งละ 1 ช่อง และไม่สามารถเดินทะลุกำแพง (ช่องที่มีค่าเป็น 1) ได้ แผนผังเขาวงกตถูกแทนด้วยตาราง Array 2 มิติ โดย \textbf{0} หมายถึงช่องที่เดินผ่านได้ และ \textbf{1} หมายถึงกำแพง

\begin{center}
\begin{tabular}{|w{c}{1cm}|*{6}{w{c}{1.8cm}|}}
\hline
& \textbf{0} & \textbf{1} & \textbf{2} & \textbf{3} & \textbf{4} & \textbf{5} \\
\hline
\textbf{0} & \textbf{0 (เริ่ม)} & 0 & 1 & 0 & 1 & 0 \\
\hline
\textbf{1} & 0 & 1 & 0 & 0 & 0 & 0 \\
\hline
\textbf{2} & 0 & 1 & 0 & 1 & 0 & 1 \\
\hline
\textbf{3} & 0 & 0 & 0 & 1 & 0 & 0 \\
\hline
\textbf{4} & 1 & 1 & 0 & 0 & 1 & 0 \\
\hline
\textbf{5} & 0 & 0 & 0 & 1 & 0 & \textbf{0 (จบ)} \\
\hline
\end{tabular}
\end{center}

จงหาว่าหุ่นยนต์จะต้องใช้จำนวน\textbf{ก้าว (Steps)} ที่สั้นที่สุดกี่ก้าว จึงจะเดินทางจากจุดเริ่มต้นไปยังจุดสิ้นสุดได้
\newline\newline
ตอบ ในกระดาษคำตอบ \newline
% answer: 14
% solution: BFS แผ่ขยายจาก (0,0) ทีละ Layer จนถึง (5,5) ที่ Layer 14
%   L0 : (0,0)
%   L1 : (0,1), (1,0)
%   L2 : (2,0)
%   L3 : (3,0)
%   L4 : (3,1)
%   L5 : (3,2)
%   L6 : (4,2), (2,2)           ← ทางแยกแรก
%   L7 : (4,3), (5,2), (1,2)
%   L8 : (5,1), (1,3)            ← (4,3) ตัน
%   L9 : (5,0), (1,4), (0,3)    ← (5,0) ตัน, (0,3) ตัน
%   L10: (1,5), (2,4)
%   L11: (0,5), (3,4)            ← (0,5) ตัน
%   L12: (3,5)
%   L13: (4,5)
%   L14: (5,5) ✓                 ← เจอจุดหมาย!
%   เส้นทางสั้นสุด 14 ก้าว: (0,0)→(1,0)→(2,0)→(3,0)→(3,1)→(3,2)→(2,2)→(1,2)→(1,3)→(1,4)→(2,4)→(3,4)→(3,5)→(4,5)→(5,5)
